\naslov{Slučajne spremenljivke in porazdelitve}

\begin{definicija}
  \pojem{Slučajna spremenljivka} $X$ je funkcija $\Omega \to \R$, da je za $a <
  b$ množica $X^{-1}((a, b])$ dogodek v $\sigma$-algebri dogodkov $\mathcal{F}$.
\end{definicija}

\begin{opomba}
  Ekvivalentno definicijo dobimo, če namesto polodprtih intervalov vzamemo
  odprte ali zaprte.
  Izbiro je predpisal ISO standard.
\end{opomba}

\begin{opomba}
  Funkcija sama po sebi je popolnoma deterministična, naključna je izbira
  argumenta.
\end{opomba}

\begin{definicija}
  Slučajna spremenljivka je \pojem{diskretna}, če je njena zaloga vrednosti
  števna ali končna množica.
\end{definicija}

\begin{definicija}
  \pojem{Porazdelitev} diskretne slučajne spremenljivke $X$ z vrednostmi
  $(x_i)_i$ je dana z verjetnostmi $P(X^{-1}(x_i))$.
\end{definicija}

\vprasanje{Definiraj slučajne spremenljivke. Kdaj je slučajna spremenljivka
  diskretna? Kaj je porazdelitev?}

Obstaja nekaj standardnih diskretnih porazdelitev.

\begin{primer}[Hipergeometrijska porazdelitev]
  Imamo posodo z $B$ belimi in $R$ rdečimi kroglicami.
  Označimo $N = B+R$ in naključno izberemo $n \le N$ kroglic tako, da so vse
  podmnožice enako verjetne.
  Če z $X$ označimo število izbranih belih kroglic, dobimo slučajno
  spremenljivko.
  Za $\max\{0, n-R\} \le k \le \min\{n, B\}$ je
  \[
	P(X = k) = \frac{\binom{B}{k} \binom{R}{n-k}}{\binom{N}{n}}.
  \]
  Na kratko označimo $X \sim \operatorname{HiperGeom}(n, B, N)$.
\end{primer}

\vprasanje{Opiši hipergeometrijsko porazdelitev.}

\begin{primer}[Binomska porazdelitev]
  Kovanec z maso $m$ vržemo $n$-krat zaporedoma, pri čemer so vsi meti
  medsebojno neodvisni, verjetnost grba pa je $p \in (0,1)$.
  Naj bo $X$ število grbov v teh $n$ metih.
  Tedaj za $k = 0, \ldots, n$ velja
  \[
	P(X = k) = \binom{n}{k} p^k (1-p)^{n-k}.
  \]
  Označimo $X \sim \operatorname{Bin}(n, p)$.
\end{primer}

\vprasanje{Opiši binomsko porazdelitev.}

\begin{primer}[Geometrijska porazdelitev]
  Naj bo $X$ število metov kovanca, potrebnih, da pade prvi grb.
  Pri tem so meti neodvisni, kovanec pade na grb z verjetnostjo $p$.
  Možne vrednosti za $X$ so vsi $k \in \N$, velja
  \[
	P(X = k) = (1-p)^{k-1} p.
  \]
  Na kratko označimo $X \sim \operatorname{Geom}(p)$.
\end{primer}

\vprasanje{Opiši geometrijsko porazdelitev.}

\begin{primer}[Negativna binomska porazdelitev]
  Mečemo kovanec in čakamo na $m$ grbov; naj bo $X$ število potrebnih metov.
  Možne vrednosti $X$ so tedaj $k = m, m+1, \ldots$, pri čemer velja
  \[
	P(X=k) = \binom{k-1}{m-1} p^m (1-p)^{k-m}.
  \]
  Oznaka je $X \sim \operatorname{NegBin}(m, p)$.
\end{primer}

\vprasanje{Opiši negativno binomsko porazdelitev.}

\begin{definicija}
  \pojem{Pochhammerjev simbol} $(a)_n$ je definiran kot
  \[
	(a)_n = a (a+1) \ldots (a+n-1).
  \]
\end{definicija}

\begin{opomba}
  Izračunamo ga lahko tudi kot
  \[
	(a)_n = \frac{\Gamma(a+n)}{\Gamma(a)}.
  \]
\end{opomba}

\begin{primer}[Poissonova porazdelitev]
  Oglejmo si dogajanje binomske porazdelitve, ko velja $np = \lambda$ konstanta,
  in ko $n \to \infty$.
  Tedaj
  \begin{align*}
	\lim_{n \to \infty} \binom{n}{k} \left( \frac{\lambda}{n} \right)^k
	\left( 1 - \frac{\lambda}{n} \right)^{n-k}
	&= \lim_{n \to \infty} \frac{\lambda^k}{k!} \frac{n(n-1) \ldots (n-k+1)}{n^k}
	\left( 1 - \frac{\lambda}{n} \right)^n \left( 1 - \frac{\lambda}{n}
	\right)^{-k} \\
	&= \frac{\lambda^k}{k!} e^{-\lambda}.
  \end{align*}
  Če je za $k = 0, 1, \ldots$ verjetnost
  \[
	P(X = k) = e^{-\lambda} \frac{\lambda^k}{k!},
  \]
  pravimo, da ima $X$ Poissonovo porazdelitev, in označimo $X \sim
  \operatorname{Po}(\lambda)$.
\end{primer}

\vprasanje{Opiši Poissonovo porazdelitev.}

\begin{definicija}
  Porazdelitev zvezne slučajne spremenljivke je podana z verjetnostmi
  $P(X \in (a, b])$ za $a < b$.
\end{definicija}

\begin{opomba}
  Pogosto želimo izračunati $P(X \in A)$, kjer $A \subseteq \R$ ni interval.
  V tem primeru lahko verjetnost izračunamo, če je $A$ sestavljena iz števnih
  unij, števnih presekov in komplementov polodprtih intervalov.
  Taki družini pravimo \pojem{Borelove množice}, in jo označimo z $B(\R)$.
  Tehnično so to najmanjša $\sigma$-algebra na $\R$, ki vsebuje vse polodprte
  intervale.
\end{opomba}

\begin{definicija}
  Slučajna spremenljivka $X$ ima \pojem{zvezno porazdelitev}, če obstaja
  nenegativna funkcija $f_X : \R \to [0, \infty)$, da je
  \[
	P(X \in (a,b]) = \int_a^b f_X(x) dx.
  \]
  Funkciji $f_X$ pravimo \pojem{gostota porazdelitve}.
\end{definicija}

\vprasanje{Definiraj zvezno porazdelitev in gostoto porazdelitve.}

\begin{primer}[Normalna porazdelitev]
  Pravimo, da ima $X$ normalno porazdelitev s parametroma $\mu$ in $\sigma^2$,
  če je gostota enaka
  \[
	f_X(x) = \inv{\sqrt{2\pi} \sigma} \exp\left(-\frac{(x-\mu)^2}{2 \sigma^2}\right).
  \]
  Pri tem je $\sigma$ razdalja od $\mu$ do prevoja, $\mu$ pa središče
  porazdelitve.
\end{primer}

\vprasanje{Opiši normalno porazdelitev.}

\begin{primer}[Eksponentna porazdelitev]
  Pravimo, da ima $X$ eksponentno porazdelitev s parametrom $\lambda$, če velja
  \[
	f_X(x) =
	\begin{cases}
	  \lambda e^{-\lambda x} & x \ge 0, \\
	  0 & x < 0.
	\end{cases}
  \]
  Označimo z $X \sim \operatorname{exp}(\lambda)$.
\end{primer}

\begin{primer}[Gama porazdelitev]
  Slučajna spremenljivka $X$ ima gama porazdelitev s parametroma $a, \lambda >
  0$, če je gostota enaka
  \[
	f_X(x) =
	\begin{cases}
	  \frac{\lambda^a}{\Gamma(a)} x^{a-1} e^{-\lambda x} & x \ge 0, \\
	  0 & x < 0.
	\end{cases}
  \]
  Oznaka: $X \sim \Gamma(a, \lambda)$.
\end{primer}

\begin{primer}[Enakomerna porazdelitev]
  Predvidljivo je
  \[
	f_X(x) =
	\begin{cases}
	  \inv{b-a} & a \le x \le b, \\
	  0 & \text{sicer.}
	\end{cases}
  \]
  Oznaka: $X \sim \operatorname{U}(a,b)$.
\end{primer}

\begin{primer}[Beta porazdelitev]
  Spremenljivka $X$ ima beta porazdelitev, če je
  \[
	f_X(x) =
	\begin{cases}
	  \inv{B(a, b)} x^{a-1} (1-x)^{b-1} & 0 < x < 1, \\
	  0 & \text{sicer.}
	\end{cases}
  \]
  Oznaka: $X \sim \operatorname{Beta}(a, b)$.
\end{primer}

\vprasanje{Opiši eksponentno, gama, enakomerno in beta porazdelitev.}

Če je $X$ slučajna spremenljivka, kakšna mora biti funkcija $f: \R \to \R$, da
bo $Y = f(X)$ tudi slučajna spremenljivka?
Za poljubna $a < b$ mora biti $X^{-1}(f^{-1}((a, b)))$ dogodek.
Če je funkcija (odsekoma) zvezna, že zadošča; potreben in zadosten pogoj pa je,
da je $f$ merljiva, torej da je $f^{-1}(A)$ dogodek za vse $A \in B(\R)$.

\begin{definicija}
  \pojem{Porazdelitvena funkcija} slučajne spremenljivke $X$ je funkcija $F_X :
  \R \to \R$, podana z $F_X(x) = P(X \le x)$.
\end{definicija}
Če ima $X$ gostoto $f_X$, je
\[
  F_X(x) = \int_{-\infty}^x f_X(u) du.
\]

\begin{izrek}
  Naj bo $F_X$ porazdelitvena funkcija slučajne spremenljivke $X$.
  Tedaj velja
  \begin{itemize}
  \item $F_X$ je nepadajoča,
  \item $\lim_{x \to \infty} F_X(x) = 1$ in $\lim_{x \to -\infty} F_X(x) = 0$,
  \item $F_X$ je desno zvezna.
  \end{itemize}
\end{izrek}

\begin{proof}
  Prva točka:
  za $x < y$ velja $F_X(y) - F_X(x) = P(X \in (x, y]) \ge 0$.

  Druga točka:
  Definiramo $A_n = \{X \le n\}$.
  Tedaj velja
  \[
	\bigcup_{n=1}^\infty A_n = \Omega,
  \]
  in ti dogodki so naraščajoči.
  Potem je
  \[
	1 = P(\Omega) = P\left( \bigcup_n A_n \right) = \lim_{n \to \infty} P(A_n)
	= \lim_{n \to \infty} F_x(n).
  \]
  Ker je $F_X$ nepadajoča, velja tudi $\lim_{x \to \infty} F_X(x) = 1$ zvezno.
  Za drugo formulo podobno definiramo $B_n = \{X \le -n\}$, kar je padajoče
  zaporedje dogodkov s praznim presekom.
  Za limito velja podoben sklep kot prej.

  Tretja točka:
  Naj bo $x_n \downarrow x$.
  Definiramo $C_n = \{X \le x_n\}$, velja
  \[
	\bigcap_{n=1}^{\infty} C_n = \{X \le x\}.
  \]
  Ker so $C_n$ padajoči, je
  \[
	F_X(x) = P\left( \bigcap_{n=1}^\infty C_n \right) = \lim_{n \to \infty} P(X
	\le x_n) = \lim_{n \to \infty} F_X(x_n),
  \]
  torej je $F_X$ res desno zvezna.
\end{proof}

\vprasanje{Definiraj porazdelitveno funkcijo slučajne spremenljivke. Kakšne
  lastnosti ima?}

\begin{vo}{Naj velja $X \sim N(\mu, \sigma^2)$ in $Y = aX+b$. Kakšna je gostota
	$Y$?}
  Računamo
  \[
	F_Y(y) = P(Y \le y)
	= P(X \le \frac{y-b}{a})
	= \inv{\sqrt{2\pi} \sigma} \int_{-\infty}^{(y-b)/a} \exp\left( - \frac{(u -
		\mu)^2}{2 \sigma^2} \right) du.
  \]
  Ker je $F_X$ zvezno odvedljiva, je
  \[
	f_Y(y) = F_Y'(y) = \inv{\sqrt{2 \pi} \sigma a} \exp\left( - \frac{(y-b-a
		\mu)^2}{2 a^2 \sigma^2} \right),
  \]
  torej $Y \sim N(a \mu + b, a^2 \sigma^2)$.
\end{vo}

\begin{definicija}
  Če je $Z \sim N(0, 1)$, rečemo, da ima $Z$ \pojem{standardizirano normalno
	porazdelitev}.
  Porazdelitveno funkcijo $Z$ označimo s $\phi$, torej
  \[
	\phi(z) = \inv{\sqrt{2 \pi}} \int_{-\infty}^z \exp\left( -\pol u^2 \right) du.
  \]
\end{definicija}

Če je $X \sim N(\mu, \sigma^2)$, je torej
\[
  F_X(x) = \phi\left(\frac{x-\mu}{\sigma}\right).
\]

\vprasanje{Kaj je standardizirana normalna porazdelitev?}

\begin{vo}{Kaj je verjetnostna transformacija?}
  Naj bo $X$ zvezno porazdeljena s porazdelitveno funkcijo $F_X$, za katero
  predpostavimo, da je zvezna.
  Definiramo $Y = F_X(X)$ in računamo za $y \in (0,1)$
  \[
	P(Y \le y)
	= P(X \in F^{-1}((-\infty, y]))
	= F_X(\sup\{x \such F_X(x) \le y\})
	= y,
  \]
  torej $Y \sim U(0,1)$.
\end{vo}

\begin{definicija}
  Naj bo $p \in (0,1)$.
  Vsakemu številu $x_p$, za katerega je $P(X \le x_p) = p$, rečemo \pojem{$p$-ti
	kvantil} porazdelitve slučajne spremenljivke $X$.
\end{definicija}

\begin{opomba}
  $p$-ti kvantil \poudarjenanegacija{ni} enolično določen.
\end{opomba}

\vprasanje{Kaj je $p$-ti kvantil porazdelitve slučajne spremenljivke?}